\section{Variabel: var, let, dan const}

Variabel adalah wadah untuk menyimpan nilai data. JavaScript menyediakan tiga cara mendeklarasikan variabel: \code{var} (cara lama), \code{let} (ES6, untuk nilai yang berubah), dan \code{const} (ES6, untuk nilai yang tidak berubah). Memahami perbedaan ketiganya---terutama dalam hal \textbf{scope} dan \textbf{hoisting}---adalah fundamental untuk menulis JavaScript yang benar \cite{jsInfo}.

\begin{table}[htbp]
\centering
\begin{tabular}{|P{2cm}|P{2.5cm}|P{2.5cm}|P{2cm}|P{3cm}|}
\hline
\textbf{Fitur} & \textbf{var} & \textbf{let} & \textbf{const} & \textbf{Catatan} \\
\hline
Scope & Function & Block & Block & Block = \code{\{ \}} \\
\hline
Hoisting & Ya (undefined) & Ya (TDZ) & Ya (TDZ) & TDZ = temporal dead zone \\
\hline
Re-assign & Ya & Ya & Tidak & const wajib inisialisasi \\
\hline
Re-declare & Ya & Tidak & Tidak & var bisa duplikat \\
\hline
\end{tabular}
\caption{Perbandingan var, let, dan const}
\end{table}

\begin{webcode}[language=JavaScript, caption={Perbedaan var, let, dan const}]
// var: function-scoped, bisa re-declare
var x = 10;
var x = 20;  // OK (re-declare)
console.log(x); // 20

// let: block-scoped, bisa re-assign
let y = 10;
y = 20;      // OK (re-assign)
// let y = 30; // Error! (re-declare)

// const: block-scoped, tidak bisa re-assign
const PI = 3.14159;
// PI = 3;   // Error! (re-assign)

// const dengan objek/array: isi bisa berubah
const arr = [1, 2, 3];
arr.push(4); // OK (mutasi isi, bukan referensi)
\end{webcode}

\begin{konsep}
\textbf{Rekomendasi: Gunakan \code{const} sebagai default.} Mulai dengan \code{const} untuk setiap variabel. Jika nanti perlu mengubah nilainya, ubah ke \code{let}. Hindari \code{var}---scope-nya terlalu luas (function scope) dan perilaku hoisting-nya sering menyebabkan bug. Ini bukan soal preferensi; ini tentang kode yang lebih aman dan lebih mudah dibaca \cite{mdnJsGuide}.
\end{konsep}

